% 分野別類題: 階数低下法 解答例
% コンパイル: TEXINPUTS="../../../00_common:$TEXINPUTS" latexmk -xelatex answers.tex
\documentclass[a4paper,11pt]{article}
\usepackage{preamble}
\usepackage{shoakubox}

\begin{document}

\kamokumei{類題演習 解答例 --- 階数低下法}

\daimon{【解法】階数低下法の解き方と導出}

\noindent
2階常微分方程式は、一般には解くのが難しいが、次の3つの典型的な場合には1階の方程式に帰着させて解くことができる（階数低下法）。

\medskip
\noindent
\textbf{(i) $y$ を含まない型 $F(x, y', y'') = 0$}

\noindent
$p = y'$ とおくと $p' = y''$ であるから、方程式は
\[
F(x, p, p') = 0
\]
という $p$ に関する1階の方程式になる。これを解いて $p = p(x)$ を求めた後、
\[
y = \int p(x)\,dx + C_{2}
\]
により $y$ を求める。

\medskip
\noindent
\textbf{(ii) $x$ を含まない型 $F(y, y', y'') = 0$}

\noindent
$p = y'$ とおき、$y$ を独立変数とみなして合成関数の微分を用いると
\[
y'' = \frac{dp}{dx} = \frac{dp}{dy}\cdot\frac{dy}{dx} = p\,\frac{dp}{dy}
\]
であるから、方程式は
\[
F\left(y,\ p,\ p\frac{dp}{dy}\right) = 0
\]
という $p$ と $y$ に関する1階の方程式になる。これを解いて $p = p(y)$（すなわち $\dfrac{dy}{dx} = p(y)$）を求めた後、変数分離して
\[
\int \frac{dy}{p(y)} = x + C_{2}
\]
により $y$ を（陰関数として、あるいは可能なら陽に）求める。

\medskip
\noindent
\textbf{(iii) 既知の1解から第2解を求める型}

\noindent
2階線形同次方程式
\[
y'' + P(x)\,y' + Q(x)\,y = 0
\]
の1つの解 $y_{1}$（$y_{1} \not\equiv 0$）が既知のとき、$y = y_{1}(x)\,v(x)$ とおいて代入すると、$y_{1}$ が解であることから $v$ を含まない項（$v$ の係数の項）が消え、$v$ に関する方程式は $v''$ と $v'$ のみを含む形になる。すなわち $w = v'$ とおけば $w$ に関する1階線形方程式に帰着し、階数が1つ下がる。これを解いて $v$ を求めれば、
\[
y_{2} = y_{1}\,v
\]
が $y_{1}$ と1次独立な第2の解となり、一般解は $y = C_{1} y_{1} + C_{2} y_{2}$ で与えられる。

\clearpage
\begin{mondaiboxS}{問1}
$x\,y'' = y'$ の一般解を求めよ。
\end{mondaiboxS}
\kaitou
$y$ を含まない型なので $p = y'$ とおくと $p' = y''$ であるから
\[
x\,p' = p
\]
$p \neq 0$ のとき変数分離して
\[
\frac{dp}{p} = \frac{dx}{x}
\quad\Longrightarrow\quad
\ln|p| = \ln|x| + C_{1}
\quad\Longrightarrow\quad
p = C_{1}' x
\]
（$p \equiv 0$ すなわち $y$ が定数の場合も $C_{1}' = 0$ として含まれる。）よって $y' = C_{1}' x$ であるから、積分して
\[
y = \frac{C_{1}'}{2}x^{2} + C_{2}
\]
定数を $A = \dfrac{C_{1}'}{2}$ とおき直せば
\[
\boxed{\; y = A\,x^{2} + B \;}
\qquad (A,\ B \text{ は任意定数})
\]
（検算: $y' = 2Ax$, $y'' = 2A$ より $x\,y'' = 2Ax = y'$。成立。）

\clearpage
\begin{mondaiboxS}{問2}
$y\,y'' = (y')^{2}, \quad y(0) = 1,\ y'(0) = 2$ を解け。
\end{mondaiboxS}
\kaitou
$x$ を含まない型なので $p = y'$, $y'' = p\dfrac{dp}{dy}$ とおくと
\[
y\,p\,\frac{dp}{dy} = p^{2}
\]
$p \neq 0$（$y'(0) = 2 \neq 0$ よりこの範囲で妥当）のとき両辺を $p$ で割って
\[
y\,\frac{dp}{dy} = p
\quad\Longrightarrow\quad
\frac{dp}{p} = \frac{dy}{y}
\quad\Longrightarrow\quad
\ln|p| = \ln|y| + C_{1}
\quad\Longrightarrow\quad
p = C_{1}'\,y
\]
すなわち $\dfrac{dy}{dx} = C_{1}'\,y$。再び変数分離して
\[
\frac{dy}{y} = C_{1}'\,dx
\quad\Longrightarrow\quad
y = C_{2}\,e^{C_{1}' x}
\]
初期条件 $y(0) = 1$ より $C_{2} = 1$。また $y' = C_{1}' C_{2} e^{C_{1}' x}$ に $x=0$ を代入して $y'(0) = C_{1}' C_{2} = C_{1}' = 2$。よって
\[
\boxed{\; y = e^{2x} \;}
\]
（検算: $y' = 2e^{2x}$, $y'' = 4e^{2x}$ より $y\,y'' = e^{2x}\cdot 4e^{2x} = 4e^{4x} = (2e^{2x})^{2} = (y')^{2}$。また $y(0)=1$, $y'(0)=2$ も満たす。成立。）

\clearpage
\begin{mondaiboxS}{問3}
$x^{2} y'' - 2x\,y' + 2y = 0 \ (x \neq 0)$ は $y_{1} = x$ を解にもつ。第2の解 $y_{2}$ と一般解を求めよ。
\end{mondaiboxS}
\kaitou
$y_{1} = x$ が解であることは、$y_{1}' = 1$, $y_{1}'' = 0$ を代入して
$x^{2}\cdot 0 - 2x\cdot 1 + 2x = -2x + 2x = 0$
より確認できる。

$y = x\,v(x)$ とおくと
\[
y' = v + x v', \qquad y'' = 2v' + x v''
\]
これを方程式に代入すると
\[
x^{2}(2v' + x v'') - 2x(v + x v') + 2x v = 0
\]
\[
2x^{2} v' + x^{3} v'' - 2x v - 2x^{2} v' + 2x v = 0
\]
$-2x^2 v'$ と $2x^2 v'$ が打ち消し合い、$-2xv$ と $2xv$ も打ち消し合うので
\[
x^{3} v'' = 0
\quad\Longrightarrow\quad
v'' = 0 \quad (x \neq 0)
\]
よって $v' = C_{1}$、積分して $v = C_{1} x + C_{2}$。したがって
\[
y = x\,v = C_{1} x^{2} + C_{2} x
\]
$y_{1} = x$ に対応する項 $C_{2}x$ を除いた $y_{2} = x^{2}$ が第2の解であり、一般解は
\[
\boxed{\; y = C_{1} x^{2} + C_{2}\,x \;}
\qquad (C_{1},\ C_{2} \text{ は任意定数})
\]
（検算: $y_{2} = x^{2}$ について $y_{2}' = 2x$, $y_{2}'' = 2$ より $x^{2}\cdot 2 - 2x\cdot 2x + 2x^{2} = 2x^{2} - 4x^{2} + 2x^{2} = 0$。成立。$y_{1} = x$ と $y_{2} = x^{2}$ はロンスキアン $W = x\cdot 2x - 1\cdot x^{2} = x^{2} \neq 0\ (x\neq0)$ より1次独立。）

\end{document}
